\subsection{Integrated Environmental Monitoring System (SIMA)}
\label{sec:sima}

The Integrated Environmental Monitoring System (SIMA) is the network of the
Government of Nuevo León responsible for measuring and monitoring air
quality. Its purpose is to determine the pollutant concentrations to which
the population is exposed, to issue alerts when pollution is high, and to
keep citizens informed about environmental health and air quality
\parencite{gobiernoNL2026PIGECA}.

SIMA began operating on 20 November 1992 as a department of the
Undersecretariat for Environmental Protection and Natural Resources, with
five monitoring stations. Over 32 years it has expanded and modernized its
infrastructure \parencite{villasaez2024}. It currently operates 15 fixed
stations distributed across 11 municipalities of the Monterrey metropolitan
area, plus 2 mobile stations \parencite{gobiernoNL2026ModernizacionSIMA}.

SIMA measures eight criteria pollutants and seven meteorological parameters,
listed with their units and coverage in the data dictionary
(\cref{tab:dic-contaminantes,tab:dic-meteo}).

SIMA is the government infrastructure that turns information on air
pollution into environmental management actions, and it is the direct
source of the data used in this work.
